\subsection{Comparison with North American practice}
\label{sec:de-comparison}

Table~\ref{tab:de-comparison} compares the connections of
Section~\ref{sec:na} with German practice.

\begin{table}[!ht]
  \centering
  \small
  \begin{tabularx}{\linewidth}{@{}>{\raggedright\arraybackslash}p{0.22\linewidth}
                                 >{\raggedright\arraybackslash}X
                                 >{\raggedright\arraybackslash}X@{}}
    \toprule
    \multicolumn{1}{c}{\textbf{Item}} & \textbf{North America (480\,V to 240\,V)} & \textbf{Germany} \\
    \midrule
    Utilization voltages   & 480Y/277, 240, 208Y/120\,V; 120/240\,V single-phase & 400/230\,V; 690\,V in industry \\
    Frequency              & 60\,Hz & 50\,Hz \\
    Single-phase supply    & 120/240\,V split-phase, centre-tapped & 230\,V line to neutral \\
    Location of the LV transformer & inside the building, from 480\,V & MV/LV substation; LV-to-LV is the exception \\
    Delta--wye phase shift & lower-voltage side lags 30\,$^{\circ}$ (Dyn1) & Dyn5, 150\,$^{\circ}$ \\
    Terminal markings      & H1--H3, X0--X3 & 1U/1V/1W, 2U/2V/2W, 2N \\
    Unearthed systems      & ungrounded or corner-grounded delta & IT system with insulation monitoring \\
    Centre-tapped delta    & in use & not used \\
    Rules                  & NEC (NFPA~70), IEEE~C57.12 series & DIN~VDE~0100, VDE-AR-N~4100/4110, EN~50588-1, IEC~60076 \\
    \bottomrule
  \end{tabularx}
  \caption{North American and German low-voltage transformer practice.}
  \label{tab:de-comparison}
\end{table}
\FloatBarrier
